\subsection{Batched Reassociation: Program and Correctness}
\label{app:batched-reassociation-program}
\label{extended:sec:partitioning}
Consider batched reassociation:
\begin{equation}
 T=\operatorname{dot}(X,Y),\quad O=\operatorname{dot}(T,Z)
 \quad\Rightarrow\quad
 U=\operatorname{dot}(Y,Z),\quad O'=\operatorname{dot}(X,U).
\end{equation}
The following program gives all eight RHS attribute lists. Let $b_X,b_Y$ be the inner source dot's batch lists,
$b_Z$ the outer source dot's right batch list, and $b=|b_X|$. The guard
requires the outer left batch list to equal $\operatorname{range}(b)$,
the positions of the inner batch group $\vec B$, in the same order. It imposes
no fixed rank, contraction-list length, or contiguity restriction.
For checked metadata arithmetic, we additionally require
$n_X+n_Y+n_Z\le\mathsf{Max}$, where $n_W=\operatorname{rank}(W)$ and
$\mathsf{Max}=2^{63}-1$. This conservative bound ensures representable
lengths and offsets; it is an explicit guard, not an inferred condition.

Let $k_X,k_Y$ be the inner source dot's contraction lists and
$d_\ell,d_r$ the outer source dot's lists in $T,Z$. The inner free groups
are $\vec I$ from $X$ and $\vec J$ from $Y$, with shared contracted group
$\vec K$. Write $\lbrack\!\lbrack e\rbrack\!\rbrack_p$ for partial metadata
evaluation. Using the paired matching of
Section~\ref{sec:symbolic-axis-partitioning}, the following expressions
use only source shapes and attributes:
\begin{equation}
\begin{aligned}
 i_X&=\ocomp(\operatorname{range}(n_X),b_X\Vert k_X),\\
 j_Y&=\ocomp(\operatorname{range}(n_Y),b_Y\Vert k_Y),\\
 q_I&=b+\operatorname{range}(|i_X|),\\
 q_J&=b+|i_X|+\operatorname{range}(|j_Y|),\\
 (u_G,v_G)&=\mpo(d_\ell,q_G),\quad z_G=\gath(d_r,u_G)
                 \quad(\vec G\in\{\vec I,\vec J\}),\\
 f_Y&=\ocomp(\operatorname{range}(n_Y),
                      b_Y\Vert\gath(j_Y,v_J)),\\
 f_Z&=\ocomp(\operatorname{range}(n_Z),b_Z\Vert z_J),\\
 s_I&=b+|f_Y|+\pos(f_Z,z_I),\qquad
 s_K=b+\pos(f_Y,k_Y).
\end{aligned}
\label{eq:batched-program-locals}
\end{equation}
Here $u_G$ indexes slots in the outer contraction list, while $v_G$
indexes axes within the producer group; only $u_G$ selects the partners
in $d_r$. The outer contraction group is
$\vec D=\gath(\operatorname{Ord}_T,d_\ell)$.
For $\vec G\in\{\vec I,\vec J\}$, define its selected subgroup as
$\vec G_{\mathrm{sel}}=\gath(\vec G,v_G)$; the remaining subgroup
$\vec G_{\mathrm{rem}}$ preserves the other axes in their original order.
Here $\pos(a,b)$ returns the positions of entries of $b$ in $a$, in
$b$'s order. The lists $q_I,q_J$ account for the batch prefix of the source interface
$\vec B\Vert \vec I\Vert \vec J$. The final two expressions locate the selected $\vec I$
axes and $\vec K$ in the replacement interface
\begin{equation}
 \operatorname{Ord}_U=\vec B\Vert\gath(\operatorname{Ord}_Y,f_Y)
                  \Vert\gath(\operatorname{Ord}_Z,f_Z).
 \label{eq:batched-u-interface}
\end{equation}
The source-derived lengths supply every offset; the executable program
contains no query of an unknown RHS interface.

\begin{table}[!htbp]
\caption{Complete metadata program for guarded batched reassociation.
Locals are defined in Equation~\eqref{eq:batched-program-locals}.}
\label{tab:reassociation-program}
\small
\begin{tabularx}{\linewidth}{@{}lYY@{}}\toprule
Attribute & $U=\operatorname{dot}(Y,Z)$ & $O'=\operatorname{dot}(X,U)$\\\midrule
Left batch & $b_Y$ & $b_X$\\
Right batch & $b_Z$ & $\operatorname{range}(b)$\\
Left contraction & $\gath(j_Y,v_J)$ & $\gath(i_X,v_I)\Vert k_X$\\
Right contraction & $z_J$ & $s_I\Vert s_K$\\\bottomrule
\end{tabularx}
\end{table}

\subsubsection{Correctness}
\label{app:reassociation-validation}
Let $C$ be the batch-preservation and integer-bound guard above.
\begin{proposition}[Guarded reassociation program]
\label{prop:reassociation-program}
For every legal source match satisfying $C$, the program in
Table~\ref{tab:reassociation-program} evaluates to legal attributes and
preserves the ordered output shape and values under exact real arithmetic.
Ranks, batch lengths, and contraction lengths may vary; free groups may
be interleaved in the operands. Zero and singleton extents are allowed.
\end{proposition}
\label{app:reassociation-proof}
\begin{proof}[Proof of Proposition~\ref{prop:reassociation-program}]
\emph{Recover the source pairings.}
Source legality makes $b_X,k_X,i_X$ a partition of the positions of $X$,
and similarly for $Y$. The guard aligns the same ordered batch group $\vec B$
in all three inputs. Since the outer contraction is disjoint from that
batch prefix, every entry of $d_\ell$ lies in exactly one of $q_I,q_J$.
Thus $u_I,u_J$ partition the slots of $d_\ell$. Selecting the same slots
in its injective partner list $d_r$ gives disjoint $z_I,z_J$ covering
$d_r$. Matching and gather composition establish the actual coordinate
pairings $\Isel,\Jsel$; equal extents alone are never used to identify axes.

\emph{Prove definedness and legality.}
The critical lookups are $\pos(f_Z,z_I)$ and $\pos(f_Y,k_Y)$.
Every entry of $z_I$ is disjoint from $b_Z$ by source legality and from
$z_J$ by the partition above, so it survives in $f_Z$.
Likewise, $k_Y$ is disjoint from $b_Y$ and from the selected free positions
$\gath(j_Y,v_J)$, so it survives in $f_Y$. Distinctness and inclusion
prove both lookups defined. The shifts place $s_I$ in the right-free block
of $U$ and $s_K$ in its left-free block. These blocks are disjoint from
each other and from the batch prefix. The left contraction list also
joins disjoint selections from $i_X$ and $k_X$. Consequently all four
contraction lists are injective, in range, and disjoint from their batch
lists. Both operands select $\Jsel$ in the inner replacement and
$\Isel\Vert \vec K$ in the outer replacement, with matching extents inherited
from the source pairings. Batch extents agree for the same reason.
All list lengths and shifted positions are bounded by $n_X+n_Y+n_Z$,
so the written integer bound establishes representability as well.
Both replacement dots are therefore legal.

\emph{Check ordered outputs and reads.}
Removing $\vec K$ from the left-free block of $U$ leaves $\Jrem$ in $Y$'s
order; removing $\Isel$ from its right-free block leaves the source
outer-free group $\vec H$ in $Z$'s order. The final left-free block is $\Irem$
in $X$'s order. Selection/complement transport therefore gives
\begin{equation}
 \operatorname{Ord}_{O'}=\vec B\Vert\Irem\Vert\Jrem\Vert \vec H=\operatorname{Ord}_O.
 \label{eq:reassociation-output-proof}
\end{equation}
The established pairings transport both sides' reads into a common
coordinate frame. Let $t,i,j,h$ be coordinates of
$\vec B,\Irem,\Jrem,\vec H$, and let $a,c,k$ range over
$\Isel,\Jsel,\vec K$. Write $\widehat X,\widehat Y,\widehat Z$ for input
reads reconstructed in their original physical orders using the source
position maps; this notation does not assume contiguous operand groups.
Compositional lowering gives
\begin{equation}
\begin{aligned}
 L[t,i,j,h]
 &=\sum_{a,c}\left(\sum_k
       \widehat X(t,i,a,k)\widehat Y(t,j,c,k)\right)
       \widehat Z(t,h,a,c),\\
 R[t,i,j,h]
 &=\sum_{a,k}\widehat X(t,i,a,k)
       \left(\sum_c\widehat Y(t,j,c,k)\widehat Z(t,h,a,c)\right).
\end{aligned}
\label{eq:reassociation-lowered}
\end{equation}
The map from $\vec D$ coordinates to $(a,c)$ selects slots $u_I,u_J$;
the partition makes it a bijection with inverse placing those coordinates
back in their original slots. It preserves extents and all three reads.
Moreover, $Z$ is independent of $k$ and $X$ is independent of $c$.
The core's checked reduction-extraction and reordering laws thus normalize
both expressions to
\begin{equation}
 \sum_{a,c,k}\widehat X(t,i,a,k)
                 \widehat Y(t,j,c,k)\widehat Z(t,h,a,c).
 \label{eq:reassociation-common-normal-form}
\end{equation}
This is equality of read factors and reduction domains, including empty
ones, not only a formal use of associativity. Ordered-axis equality gives
ordered-shape equality; source dtype agreement is inherited by both new
dots. Together with definedness and legality, these establish
$\operatorname{Valid}(m,f(m))$. Lemma~\ref{lem:metadata-emission} transfers
the result to the displayed executable metadata expressions.
\end{proof}


\paragraph{Implementation evidence.}
\label{app:reassociation-record}
The displayed expressions correspond to emitted compiler output, whose
semantic check returned \code{Proved} under the written guard; compilation
and portable validation also succeeded. Exact sources, generated JSON,
commands, and provenance are in \nolinkurl{evaluation/reassociation/};
the temporary-variable correspondence is archived in
\nolinkurl{evaluation/supplementary-materials/}.
This execution record is separate from the corpus totals and is not a
checked proof certificate.

